\documentclass[12pt,a4paper]{article}
\usepackage[utf8]{inputenc}
\usepackage[spanish]{babel}
\usepackage[margin=2.5cm]{geometry}
\usepackage{natbib}
\bibliographystyle{IEEEtran}


\title{PE5 -- Vista previa individual\\Fundamento Teórico: Fundamentos de gestión de proyecto y trabajo en equipo}
\author{Nieves Sánchez, Jimmy Samuel}
\date{}

\begin{document}
\maketitle


\section*{Fundamento Teórico}

\subsection*{Fundamentos de gestión de proyecto y trabajo en equipo}

La gestión del proyecto integrador PE1--PE5 de MundiPets se organiza siguiendo los
ocho dominios de desempeño de la 7.ª edición del PMBOK Guide: partes interesadas,
equipo, enfoque de desarrollo y ciclo de vida, planificación, trabajo del proyecto,
entrega, medición e incertidumbre \cite{pmbok7}. A diferencia de ediciones anteriores
centradas en procesos secuenciales, la 7.ª edición adopta principios orientados a
valor y desempeño adaptable, lo que resulta coherente con el proceso realmente
seguido por el equipo: iteraciones cortas por práctica (PE1 a PE5), entregables
incrementales del ERS y retroalimentación continua del docente, en lugar de una
planificación cerrada de una sola vez \cite{pmbok7}.

Esta forma de trabajar exige, además, una fundamentación propia de la ingeniería de
requisitos ágil. La revisión sistemática de Hoy y Xu identifica los retos recurrentes
de conciliar la especificación rigurosa de requisitos con ciclos de entrega cortos:
comunicación insuficiente entre roles, priorización inestable y trazabilidad que se
degrada cuando el equipo no documenta explícitamente sus decisiones \cite{hoy2023}.
El marco conceptual que proponen centrado en comunicación estructurada,
priorización basada en valor y gestión activa del cambio es directamente aplicable
al reparto de responsabilidades y a las tres dependencias
críticas del cronograma: un retraso en la entrega de la
matriz de trazabilidad propaga directamente el riesgo hacia las tareas dependientes
de auditoría de métricas y de requisitos de IA, exactamente el tipo de fragilidad de
coordinación que documentan Hoy y Xu \cite{hoy2023}.

Bajo este marco, la evidencia de Git (commits distribuidos por integrante, ramas por
criterio de la rúbrica, mensajes descriptivos y frecuencia mínima semanal) no es un
simple registro técnico de cambios en el código fuente: es un artefacto de gestión en
el sentido que le da el PMBOK al dominio de desempeño de medición \cite{pmbok7},
porque permite verificar objetivamente sin depender de declaraciones subjetivas
el cumplimiento del plan de trabajo, la distribución real de la carga entre los cinco
integrantes y el factor de ajuste individual (FAI) que exige la rúbrica de la PE5.

\bibliography{nieves_referencias}

\end{document}
